% Part: proof-theory
% Chapter: natural-deduction
% Section: translation-G2i

\documentclass[../../../include/open-logic-section]{subfiles}

\begin{document}

\olfileid{pt}{nat}{tgi}

\olsection{从 \Log{G2i} 到 \Log{N2i} 的翻译}

在 \Log{G2i} 中，相继式的前件是公式的多重集~$\Gamma$；而在 \Log{N2i} 中，相继式的前件是带标签公式的集合~$\Gamma'$，且满足每个标签至多出现一次。若对每个公式~$!A$，$\Gamma'$ 中形如 $x : !A$ 的元素数目不超过~$!A$ 在~$\Gamma$ 中的重数（即~$!A$ 的副本数），则称带标签公式的集合~$\Gamma'$ \emph{对应于}多重集~$\Gamma$。

\begin{prop}\ollabel{prop:G2i-to-N2i}
如果 $\Log{G2i} + \Cut \Proves \Gamma \Sequent \Delta$，则 $\Log{N2i}
\Proves \Gamma' \Sequent !C$，其中若 $\Delta$ 为空则 $!C \ident \lfalse$，否则 $\Delta = \{!A\}$（此时 $!C \ident !A$），且 $\Gamma'$ 对应于~$\Gamma$。
\end{prop}

\begin{proof}
  我们要证明：如果 $\pi$ 是 \Log{G2ci} 中 $\Gamma \Sequent \Delta$ 的推导，则存在对应于~$\Gamma$ 的~$\Gamma'$，以及 \Log{N2i} 中~$\Gamma' \Sequent !C$ 的推导~$\delta$。为此，我们对 $n = \pheight{\pi}$ 进行归纳。

  若 $n = 0$，则 $\pi$ 是一条公理。若为公理 $\lfalse \Sequent \quad$，则 $\delta$ 为公理 $x:\lfalse \Sequent \lfalse$，即 $\Gamma' = \{x:\lfalse\}$ 且 $!C \ident \lfalse$。若 $\pi$ 是形如 $!A \Sequent !A$ 的公理，则 $\delta$ 为 $x: !A \Sequent !A$。

  若 $n>0$，则 $\pi$ 的最后一步应用了某条规则~$R$。归纳假设保证了在~\Log{N2c} 中存在对应于该规则前提的相继式的推导。我们可以假设这些推导中的所有解除标签互不相同，并且若标签 $x$ 被解除，则它只出现在解除它的推理上方（必要时可进行重标记）。接下来我们展示如何组合这些推导，以得到合适的 $\Gamma' \Sequent !C$ 的推导。我们根据~$R$ 分情况讨论。

  \begin{enumerate}
    \item 若 $R$ 为 \RightR{\Weakening}，则 $\pi$ 的最后一步为
    \[
    \AxiomC{}
    \RightLabel{$\pi_1$}
    \Deduce$\Gamma \fCenter$
    \RightLabel{\RightR{\Weakening}}
    \UnaryInf$\Gamma \fCenter !A$
    \DisplayProof
    \]
    对 $\pi_1$ 应用归纳假设，可得~$\Gamma' \Sequent \lfalse$ 的推导。我们再应用 $\FalseInt$ 规则，便得到 $\Gamma' \Sequent !A$ 的推导。
    \item 若 $R$ 为 \LeftR{\Weakening}，则 $\pi$ 的最后一步为
    \[
    \AxiomC{}
    \RightLabel{$\pi_1$}
    \Deduce$\Gamma \fCenter \Delta$
    \RightLabel{\RightR{\Weakening}}
    \UnaryInf$!A, \Gamma \fCenter \Delta$
    \DisplayProof
    \]
    对 $\pi_1$ 应用归纳假设，可得~$\Gamma' \Sequent !C$ 的推导，我们将其取作 $\delta$。注意，若 $\Gamma'$ 对应于 $\Gamma$，则它也对应于 $!A, \Gamma$，因为 $\Gamma'$ 中形如 $x : !A$ 的元素数不超过 $\Gamma$ 中~$!A$ 的副本数，从而也不超过 $\Gamma \cup \{!A\}$ 中 $!A$ 的副本数。
    \item 若 $R$ 为 \LeftR{\Contraction}，则 $\pi$ 的最后一步为
    \[
    \AxiomC{}
    \RightLabel{$\pi_1$}
    \Deduce$!A, !A, \Gamma \fCenter \Delta$
    \RightLabel{\LeftR{\Contraction}}
    \UnaryInf$!A, \Gamma \fCenter \Delta$
    \DisplayProof
    \]
    对 $\pi_1$ 应用归纳假设，可得~$\Pi, \Gamma' \Sequent !C$ 的推导，其中 $\Pi \subseteq \{x : !A,
    y:!A\}$。若 $\Pi = \{x : !A, y:!A\}$，则将~$\pi_1$ 中所有带标签公式 $y:!A$ 替换为 $x:!A$。由于 $x$ 出现在~$\pi_1$ 的末端相继式的语境中，我们可以假设~$\pi_1$ 中没有推理会解除形如~$x:!B$ 的公式，即~$\pi_1$ 中相继式左侧的 $x:!B$ 都不是主动的。因此，替换后的结果仍然是~\Log{N2i} 中的一个推导。
    \item 若 $R$ 为 $\RightR{\lif}$，则 $\pi$ 的最后一步为
    \[
    \AxiomC{}
    \RightLabel{$\pi_1$}
    \Deduce$!A, \Gamma \fCenter !B$
    \RightLabel{\RightR{\lif}}
    \UnaryInf$\Gamma \fCenter !A \lif !B$
    \DisplayProof
    \]
    由归纳假设，存在 $x: !A, \Gamma' \Sequent !B$ 或 $\Gamma' \Sequent !B$ 的推导~$\delta_1$，其中 $\Gamma'$ 对应于~$\Gamma$。我们可以在 $\delta_1$ 之后应用 \Intro{\lif} 规则，将其结果作为 $\delta$。
    \item 若 $R$ 为 \LeftR{\lif}，则 $\pi$ 的最后一步为
    \[
    \AxiomC{}
    \RightLabel{$\pi_1$}
    \Deduce$\Gamma_1 \fCenter !A$
    \AxiomC{}
    \RightLabel{$\pi_2$}
    \Deduce$!B, \Gamma_2 \fCenter \Delta$
    \RightLabel{\LeftR{\lif}}
    \BinaryInf$!A \lif !B, \Gamma_1, \Gamma_2 \fCenter \Delta$
    \DisplayProof
    \]
    由归纳假设，存在 $\Gamma_1' \Sequent !A$ 的推导 $\delta_1$，以及 $x: !B, \Gamma_2' \Sequent !C$ 或 $\Gamma_2' \Sequent !C$ 的推导 $\delta_2$。若 $\delta_2$ 为后一种形式，则可直接将 $\delta_2$ 作为~$\delta$。否则，我们可以重标记~$\pi_1$ 中的所有公式和解除标签，以确保没有标签同时出现在 $\delta_1$ 和~$\delta_2$ 中。于是 $\Gamma_1' \cup \Gamma_2'$ 对应于 $\Gamma_1 \cup \Gamma_2$。由于 $x:!B$ 出现在~$\delta_2$ 的末端相继式中，$x:!B$ 在~$\delta_2$ 的任何带解除标签~$x$ 的推理中都不可能是主动的。将~$\delta_2$ 中的每条公理 $x:!B \Sequent !B$ 替换为推导~$\delta_3$
    \[
    \Axiom$x: !A \lif !B \fCenter !A \lif !B$
    \AxiomC{}
    \RightLabel{$\delta_1$}
    \Deduce$\Gamma_1' \fCenter !A$
    \RightLabel{\Elim{\lif}}
    \BinaryInf$x: !A \lif !B \fCenter !B$
    \DisplayProof
    \]
    并将~$\delta_2$ 中每一处 $x:!B$ 替换为 $x:!A \lif !B$。
    所得结果 $\Subst{\delta_2}{\delta_3}{x:!B}$ 即为 $x:!A \lif !B, \Gamma_1', \Gamma_2' \Sequent !C$ 的推导。
    \item 若 $R$ 为 \RightR{\land}，则 $\pi$ 的最后一步为
    \[
    \AxiomC{}
    \RightLabel{$\pi_1$}
    \Deduce$\Gamma_1 \fCenter !A$
    \AxiomC{}
    \RightLabel{$\pi_2$}
    \Deduce$\Gamma_2 \fCenter !B$
    \RightLabel{\RightR{\land}}
    \BinaryInf$\Gamma_1, \Gamma_2 \fCenter !A \land !B$
    \DisplayProof
    \]
    由归纳假设，存在 $\Gamma_1' \Sequent !A$ 的推导 $\delta_1$ 和 $\Gamma_2' \Sequent !A$ 的推导 $\delta_2$。通过重标记~$\pi_2$ 中的所有公式和解除标签，我们可以确保没有标签同时出现在 $\delta_1$ 和 $\delta_2$ 中。于是 $\Gamma_1' \cup \Gamma_2'$ 对应于 $\Gamma_1 \cup \Gamma_2$。我们对 $\delta_1$ 和~$\delta_2$ 的末端相继式应用 $\Intro{\land}$ 规则，便得到推导~$\delta$。
    \item 若 $R$ 为 $\LeftR{\land}$，则 $\pi$ 的最后一步为，例如，
    \[
    \AxiomC{}
    \RightLabel{$\pi_1$}
    \Deduce$!A, \Gamma \fCenter \Delta$
    \RightLabel{\LeftR{\land}}
    \UnaryInf$!A \land !B, \Gamma \fCenter \Delta$
    \DisplayProof
    \]
    由归纳假设，存在 $x: !A, \Gamma' \Sequent !C$ 或 $\Gamma' \Sequent !C$ 的推导~$\delta_1$，其中 $\Gamma'$ 对应于~$\Gamma$。在后一种情况下，我们取 $\delta_1$ 作为 $\delta$。在前一种情况下，注意到因为 $x:!A$ 出现在末端相继式中，$x:!A$ 在~$\delta_1$ 的任何带解除标签~$x$ 的推理中都不是主动的。将每条公理 $x:!A \Sequent !A$ 替换为推导~$\delta_2$
    \[
    \Axiom$x: !A \land !B \fCenter !A \land !B$
    \RightLabel{\Elim{\land}}
    \UnaryInf$x: !A \land !B \fCenter !A$
    \DisplayProof
    \]
    并将每一处 $x:!A$ 替换为 $x:!A \land !B$，即考虑
    $\Subst{\delta_1}{\delta_2}{x:!A}$。这便是 $x:!A
    \land !B, \Gamma' \Sequent !C$ 的推导。
    \item 若 $R$ 为 \Cut，则 $\pi$ 的最后一步为
    \[
    \AxiomC{}
    \RightLabel{$\pi_1$}
    \Deduce$\Gamma_1 \fCenter !A$
    \AxiomC{}
    \RightLabel{$\pi_2$}
    \Deduce$!A, \Gamma_2 \fCenter \Delta$
    \RightLabel{\Cut}
    \BinaryInf$\Gamma_1, \Gamma_2 \fCenter \Delta$
    \DisplayProof
    \]
    由归纳假设，存在 $\Gamma_1' \Sequent !A$ 的推导 $\delta_1$ 和 $x:!A, \Gamma_2' \Sequent !C$ 或 $\Gamma_2' \Sequent !C$ 的推导 $\delta_2$。在后一种情况下，取 $\delta$ 为~$\delta_2$。否则，将 $\delta_2$ 中的每条公理 $x:!A \Sequent !A$ 替换为 $\delta_1$，从而得到 $\delta$，即 $\Subst{\delta_2}{\delta_1}{x:!A}$ 对 $\Gamma_1',
    \Gamma_2' \Sequent !C$ 的推导。
  \end{enumerate}
\end{proof}

\end{document}